\begin{thebibliography}{99}

\bibitem{Achinger}
P.~Achinger, \emph{Wild ramification and $K(\pi,1)$ spaces},
Invent. Math. \textbf{210} (2017), no.~2, 453--499.
Section~6.2, Theorem~6.2.1, records the Gabber--Fujiwara comparison
for every torsion abelian sheaf.
\href{https://arxiv.org/abs/1701.03197v2}{arXiv:1701.03197v2}
(4 May 2017);
\href{https://doi.org/10.1007/s00222-017-0733-5}{doi:10.1007/s00222-017-0733-5}.

\bibitem{AMMN}
B.~Antieau, A.~Mathew, M.~Morrow, and T.~Nikolaus,
\emph{On the Beilinson fiber square},
Duke Math. J. \textbf{171} (2022), no.~18, 3707--3806.
Propositions~2.8, 4.10, 4.12, Corollary~2.10, Theorem~2.12,
and Constructions~4.6, 4.11.
Numbered locators follow
\href{https://arxiv.org/abs/2003.12541v2}{arXiv:2003.12541v2}
(29 September 2021).

\bibitem{BerkovichEtale}
V.~G.~Berkovich, \emph{\'Etale cohomology for non-Archimedean analytic spaces},
Publications Math\'ematiques de l'IH\'ES \textbf{78} (1993), 5--161.
Theorem~6.1.1.
\href{https://www.numdam.org/article/PMIHES_1993__78__5_0.pdf}{Published article}.

\bibitem{Berthelot1997Duality}
Pierre Berthelot,
\emph{Dualit\'e de Poincar\'e et formule de K\"unneth en cohomologie rigide},
Comptes Rendus de l'Acad\'emie des Sciences de Paris,
S\'erie I, Math\'ematique \textbf{325} (1997), no.~5, 493--498.
\href{https://wstein.org/people/berthelo/publis/Poincare_Kunneth.pdf}{Author's text}.

\bibitem{BesserRegulators}
A.~Besser, \emph{Syntomic regulators and $p$-adic integration I:
rigid syntomic regulators}, Israel Journal of Mathematics
\textbf{120} (2000), part B, 291--334.
Proposition~8.6(2)--(3), its proof and equation~(8.1);
Proposition~9.9, Corollary~9.10, and Proposition~9.11.
Numbered locators refer to the author's version.
\href{https://www.math.bgu.ac.il/~bessera/reg/reg.html}{Author's original version};
\href{https://doi.org/10.1007/BF02834843}{doi:10.1007/BF02834843}.

\bibitem{BesserDeJeuAlgorithm}
A.~Besser and R.~de~Jeu, \emph{Li$^{(p)}$-service? An algorithm for computing $p$-adic polylogarithms}, Math. Comp. \textbf{77} (2008), no.~262, 1105--1134.
\href{https://doi.org/10.1090/S0025-5718-07-02027-3}{doi:10.1090/S0025-5718-07-02027-3};
\href{https://research.vu.nl/ws/portalfiles/portal/2459860/220267.pdf}{Published article}.

\bibitem{BesserDeJeuCurves}
A.~Besser and R.~de Jeu, \emph{The syntomic regulator for $K_4$ of curves}, Pacific J. Math. \textbf{260} (2012), no.~2, 305--380. Equation~(1.8), p.~307, and p.~377 record the Coleman dilogarithm normalization and functional equations.
\href{https://msp.org/pjm/2012/260-2/pjm-v260-n2-p03-s.pdf}{Published article}.

\bibitem{BesserDeJeu}
A.~Besser and R.~de Jeu, \emph{The syntomic regulator for the $K$-theory of fields}, Ann. Sci. \'Ecole Norm. Sup. (4) \textbf{36} (2003), no.~6, 867--924. Equation~(1.2) and Theorems~1.6, 1.10, 1.12.
\href{https://www.numdam.org/article/ASENS_2003_4_36_6_867_0.pdf}{Published article};
\href{https://doi.org/10.1016/j.ansens.2003.01.003}{doi:10.1016/j.ansens.2003.01.003}.

\bibitem{BGTTrace}
A.~J.~Blumberg, D.~Gepner, and G.~Tabuada,
\emph{Uniqueness of the multiplicative cyclotomic trace},
Advances in Mathematics \textbf{260} (2014), 191--232, Theorem~7.4.
\href{https://arxiv.org/abs/1103.3923v3}{arXiv:1103.3923v3}.

\bibitem{BokstedtOttosen}
M.~B\"okstedt and I.~Ottosen,
\emph{A Hochschild--Kostant--Rosenberg theorem for cyclic homology},
J. Pure Appl. Algebra \textbf{221} (2017), no.~6, 1458--1493.
Section~2, especially Definition~2.1, Remark~2.2, and Proposition~2.3.
\href{https://arxiv.org/abs/1601.07412v2}{arXiv:1601.07412v2}.

\bibitem{BondarkoIvanovChowWeights}
M.~V.~Bondarko and M.~A.~Ivanov,
\emph{On Chow weight structures for $cdh$-motives with integral coefficients},
arXiv preprint, version~2, 18 August 2015. Remark~1.3.2.
\href{https://arxiv.org/abs/1506.00631v2}{arXiv:1506.00631v2}.

\bibitem{Borel}
A.~Borel, \emph{Stable real cohomology of arithmetic groups}, Ann. Sci. \'Ecole Norm. Sup. (4) \textbf{7} (1974), no.~2, 235--272. In particular, Proposition~12.2 gives the relevant rank calculation.
\href{https://doi.org/10.24033/asens.1269}{doi:10.24033/asens.1269}.

\bibitem{BouisMotivic}
T.~Bouis, \emph{Motivic cohomology of mixed characteristic schemes},
arXiv preprint, version~3, 7 September 2026, Corollaries~5.12--5.13.
\href{https://arxiv.org/abs/2412.06635v3}{arXiv:2412.06635v3}.

\bibitem{BunkeTamme}
U.~Bunke and G.~Tamme, \emph{Multiplicative differential algebraic $K$-theory and applications}, Annals of $K$-Theory \textbf{1} (2016), no.~3, 227--258. See Remark~2.6, Theorem~2.31, and Section~6.
Numbered locators refer to the published version.
\href{https://doi.org/10.2140/akt.2016.1.227}{doi:10.2140/akt.2016.1.227};
\href{https://msp.org/akt/2016/1-3/akt-v1-n3-p01-s.pdf}{Published article}.

\bibitem{CaroDAddezio2025}
D.~Caro and M.~D'Addezio,
\emph{Injectivity failure in crystalline comparisons},
arXiv preprint, version~1, 14 November 2025.
\href{https://arxiv.org/abs/2511.11444v1}{arXiv:2511.11444v1}.

\bibitem{CartanStein1952}
H.~Cartan,
\emph{Faisceaux analytiques sur les vari\'et\'es de Stein :
d\'emonstration des th\'eor\`emes fondamentaux},
S\'eminaire Henri Cartan \textbf{4} (1951--1952),
expos\'e~19, 1--15 (1952).
\href{https://numdam.org/item/SHC_1951-1952__4__A19_0/}{Published article}.

\bibitem{CisinskiDegliseMotives}
D.-C.~Cisinski and F.~D\'eglise,
\emph{Triangulated categories of mixed motives},
Springer Monographs in Mathematics, Springer, 2019.
Introduction~B, p.~xxii; Theorem~11.1.24;
Examples~11.1.25 and~11.2.3; Corollaries~14.2.14 and~14.2.19;
Theorem~16.1.4 and Corollary~16.1.7.
\href{https://arxiv.org/abs/0912.2110v4}{arXiv:0912.2110v4}.

\bibitem{ColmezNiziol}
P.~Colmez and W.~Nizio\l, \emph{On $p$-adic comparison theorems for rigid analytic varieties, I},
M\"unster J. Math. \textbf{13} (2020), 445--507.
Theorem~1.1(4), Theorem~1.3(2), Theorem~7.9, and
Sections~3.1.1, 3.2.3, 5.1, and~7.2.4.
Numbered locators follow
\href{https://arxiv.org/abs/1905.04721v2}{arXiv:1905.04721v2}
(5 October 2019).

\bibitem{DLZ2011}
Christopher Davis, Andreas Langer, and Thomas Zink,
\emph{Overconvergent de Rham--Witt cohomology},
Annales Scientifiques de l'\'Ecole Normale Sup\'erieure (4)
\textbf{44} (2011), no.~2, 197--262.
\href{https://www.numdam.org/item/10.24033/asens.2143.pdf}{Published article}.

\bibitem{DeJeuCurves}
R.~de Jeu, \emph{Towards Regulator Formulae for the $K$-Theory of Curves over Number Fields}, Compositio Math. \textbf{124} (2000), 137--194. Theorem~2.3 restates the all-embedding complex regulator formula for the symbol map.
\href{https://doi.org/10.1023/A:1026440915009}{Published article};
\href{https://arxiv.org/abs/math/9811192}{Earlier preprint}.

\bibitem{DeJeuWedge}
R.~de Jeu, \emph{Zagier's conjecture and wedge complexes in algebraic $K$-theory}, Compositio Math. \textbf{96} (1995), no.~2, 197--247. Proposition~4.1 and Remark~5.2.
\href{https://www.numdam.org/article/CM_1995__96_2_197_0.pdf}{Published article}.

\bibitem{DennisStein}
R.~K.~Dennis and M.~R.~Stein, \emph{$K_2$ of discrete valuation rings}, Advances in Mathematics \textbf{18} (1975), no.~2, 182--238. The modern sign convention used here is specified explicitly in Section~\ref{d3:sec:explicit}.
\href{https://doi.org/10.1016/0001-8708(75)90157-7}{doi:10.1016/0001-8708(75)90157-7}.

\bibitem{EsnaultViehweg}
H.~Esnault and E.~Viehweg, \emph{Deligne--Beilinson cohomology}, in \emph{Beilinson's Conjectures on Special Values of $L$-Functions}, Perspectives in Mathematics \textbf{4}, Academic Press, 1988, 43--91. See Remark~2.13, Section~3 for products, and Section~8 for Chern classes.
\href{https://page.mi.fu-berlin.de/esnault/preprints/helene/13-deligne_beilinson.pdf}{Author's institutional copy}.

\bibitem{FeldFiniteGersten2026}
N.~Feld, \emph{Finite-coefficient Gersten injectivity fails in ramified mixed
characteristic}, preprint, 2026.
\href{https://arxiv.org/abs/2608.05005v2}{arXiv:2608.05005v2},
version~2, 24 September 2026; first submitted 5 August 2026.
Theorem~1.1 and Section~5.

\bibitem{GeisserLevineCharP}
T.~Geisser and M.~Levine, \emph{The $K$-theory of fields in
characteristic $p$}, Inventiones Mathematicae \textbf{139}
(2000), no.~3, 459--493. Theorem~8.4.
\href{https://www2.rikkyo.ac.jp/web/geisser/p-part.pdf}{Published article};
\href{https://doi.org/10.1007/s002220050014}{doi:10.1007/s002220050014}.

\bibitem{Ginot}
G.~Ginot, \emph{Formules explicites pour le caract\`ere de Chern en $K$-th\'eorie alg\'ebrique}, Annales de l'Institut Fourier \textbf{54} (2004), no.~7, 2327--2355.
\href{https://www.math.univ-paris13.fr/~ginot/papers/flecarch.pdf}{Author's preprint};
\href{https://doi.org/10.5802/aif.2081}{doi:10.5802/aif.2081}.

\bibitem{GrosseKlonne}
E.~Grosse-Kl\"onne, \emph{Rigid analytic spaces with overconvergent structure sheaf}, Journal f\"ur die reine und angewandte Mathematik \textbf{519} (2000), 73--95. Proposition~3.1 gives coherent acyclicity for dagger affinoids.
\href{https://arxiv.org/abs/1408.3329v1}{arXiv:1408.3329v1}.

\bibitem{HoyoisCircle}
M.~Hoyois, \emph{The homotopy fixed points of the circle action on
Hochschild homology}, arXiv preprint, version~2, 21 April 2018.
Theorem~2.1, Lemma~2.2, and Theorem~2.3.
\href{https://arxiv.org/abs/1506.07123v2}{arXiv:1506.07123v2}.

\bibitem{HuberKings}
A.~Huber and G.~Kings, \emph{A $p$-adic analogue of the Borel regulator and the Bloch--Kato exponential map},
J. Inst. Math. Jussieu \textbf{10} (2011), no.~1, 149--190.
Section~2 gives the Chern-compatible comparison with Besser's
syntomic cohomology and the exponential, without an unramifiedness restriction.
Numbered citations refer to
\href{https://arxiv.org/abs/math/0612611v1}{arXiv:math/0612611v1}.

\bibitem{Jannsen}
U.~Jannsen, \emph{Continuous \'etale cohomology}, Mathematische Annalen \textbf{280} (1988), 207--245. Equations~(0.2), (0.4), and Section~6.
\href{https://epub.uni-regensburg.de/26684/1/jannsen12.pdf}{Published article}.

\bibitem{KerzThesis}
M.~Kerz, \emph{Milnor $K$-theory of local rings}, doctoral thesis, Universit\"at Regensburg, 2008. Theorem~4.2.5, Proposition~4.2.8(4)--(6), and Theorem~4.3.2 give the improved theory, the ordinary-to-improved surjection, and the comparison with Quillen $K$-theory whose composite is $(n-1)!$ for general local rings.
\href{https://d-nb.info/98935329x/34}{Full thesis}.

\bibitem{LevineHigherChowRevisited}
M.~Levine,
\emph{Bloch's higher Chow groups revisited},
Ast\'erisque \textbf{226} (1994), 235--320. Theorem~4.7.
\href{https://www.numdam.org/item/AST_1994__226__235_0/}{Published article}.

\bibitem{LevineIndK3}
M.~Levine, \emph{The indecomposable $K_3$ of fields}, Ann. Sci. \'Ecole Norm. Sup. (4) \textbf{22} (1989), no.~2, 255--344. Theorem~4.12 is the finite-coefficient Chern comparison used here; see also Theorem~4.13 and Section~5.1.
\href{https://doi.org/10.24033/asens.1585}{doi:10.24033/asens.1585}.

\bibitem{MerkurjevTorsion}
A.~S.~Merkurjev, \emph{On the torsion in $K_2$ of local fields}, Annals of Mathematics \textbf{118} (1983), no.~2, 375--381.
\href{https://annals.math.princeton.edu/1983/118-2/p06}{Publisher's page};
\href{https://doi.org/10.2307/2007033}{doi:10.2307/2007033}.

\bibitem{MilneDuality}
J.~S.~Milne, \emph{Arithmetic Duality Theorems}, second edition, BookSurge, 2006. Chapter~I, Corollary~2.3 and the discussion preceding Theorem~2.8 give local Galois finiteness.
\href{https://www.jmilne.org/math/Books/ADTnot.pdf}{Author's edition}.

\bibitem{MilneEtale}
J.~S.~Milne, \emph{Lectures on \'Etale Cohomology}, author's lecture notes, version~2.21, 22 March 2013. Theorem~19.1(b) gives geometric finite-coefficient finiteness.
\href{https://www.jmilne.org/math/CourseNotes/LEC.pdf}{Author's notes}.

\bibitem{MochizukiParameters}
S.~Mochizuki, \emph{Local Gersten's conjecture for regular system of parameters}, preprint, version~8, 6 August 2020. Section~\ref{d3:sec:mochizukisquare} examines Lemma~13(III) in this version.
\href{https://arxiv.org/abs/1503.07966v8}{arXiv:1503.07966v8}.

\bibitem{NikolausScholze}
T.~Nikolaus and P.~Scholze, \emph{On topological cyclic homology},
Acta Mathematica \textbf{221} (2018), no.~2, 203--409, Corollary~I.4.3.
\href{https://arxiv.org/abs/1707.01799v2}{arXiv:1707.01799v2}.

\bibitem{Petrequin2003}
D.~Petrequin, \emph{Classes de Chern et classes de cycles en
cohomologie rigide}, Bulletin de la Soci\'et\'e Math\'ematique
de France \textbf{131} (2003), no.~1, 59--121.
\href{https://www.numdam.org/article/BSMF_2003__131_1_59_0.pdf}
{Published article};
\href{https://doi.org/10.24033/bsmf.2437}{doi:10.24033/bsmf.2437}.

\bibitem{Quillen}
D.~Quillen, \emph{Higher algebraic $K$-theory I}, in \emph{Algebraic $K$-Theory I: Higher $K$-Theories}, Lecture Notes in Mathematics \textbf{341}, Springer, 1973, 85--147.
\href{https://doi.org/10.1007/BFb0067053}{doi:10.1007/BFb0067053}.

\bibitem{QuillenFinite}
D.~Quillen, \emph{On the cohomology and \(K\)-theory of the general
linear groups over a finite field}, Ann. of Math. (2)
\textbf{96} (1972), no.~3, 552--586.
\href{https://doi.org/10.2307/1970825}{doi:10.2307/1970825}.

\bibitem{SaitoCompletion}
T.~Saito, \emph{Graded quotients of ramification groups of local fields with imperfect residue fields}, Amer. J. Math. \textbf{145} (2023), no.~5, 1389--1464, Lemma~4.2.1.
\href{https://preprint.press.jhu.edu/ajm/sites/default/files/AJM-saito.pdf}{Author's version};
\href{https://arxiv.org/abs/2004.03768v2}{arXiv:2004.03768v2};
\href{https://doi.org/10.1353/ajm.2023.a907702}{doi:10.1353/ajm.2023.a907702}.

\bibitem{SilvermanAEC2009}
Joseph H. Silverman,
\emph{The Arithmetic of Elliptic Curves},
second edition, Graduate Texts in Mathematics, vol.~106,
Springer, New York, 2009.
\href{https://doi.org/10.1007/978-0-387-09494-6}{doi:10.1007/978-0-387-09494-6}.

\bibitem{StacksBlowup}
The Stacks Project Authors, \emph{The Stacks Project}, Section~31.33, ``Blowing up,'' especially the standard chart construction and tautological ideal identity in Lemmas~31.33.2 and~31.33.4.
\href{https://stacks.math.columbia.edu/tag/01OF}{Tag 01OF};
\href{https://stacks.math.columbia.edu/tag/0804}{Tag 0804};
\href{https://stacks.math.columbia.edu/tag/02OS}{Tag 02OS}.

\bibitem{StacksFlatness}
The Stacks Project Authors, \emph{The Stacks Project}, Lemma~10.128.1 (Miracle flatness), Tag~00R4.
\href{https://stacks.math.columbia.edu/tag/00R4}{Tag 00R4}.

\bibitem{StacksHenselization}
The Stacks Project Authors, \emph{The Stacks Project}, Section~10.156, Henselization and quasi-finite ring maps (Tag~0GIN), and Section~15.46, Permanence of properties under henselization (Tag~07QL).
\href{https://stacks.math.columbia.edu/tag/0GIN}{Tag 0GIN}; \href{https://stacks.math.columbia.edu/tag/07QL}{Tag 07QL}. Also Section~10.153, \href{https://stacks.math.columbia.edu/tag/04GK}{Tag 04GK}, and Section~10.155, \href{https://stacks.math.columbia.edu/tag/0BSK}{Tag 0BSK}.
For the cited permanence and quotient results, see
\href{https://stacks.math.columbia.edu/tag/07QM}{Tag 07QM}
(Lemma~15.46.1),
\href{https://stacks.math.columbia.edu/tag/06LJ}{Tag 06LJ}
(Lemma~15.46.3),
\href{https://stacks.math.columbia.edu/tag/06LN}{Tag 06LN},
\href{https://stacks.math.columbia.edu/tag/0AP3}{Tag 0AP3}
(Lemma~15.46.11),
\href{https://stacks.math.columbia.edu/tag/05WQ}{Tag 05WQ}
(Lemma~10.156.2), and
\href{https://stacks.math.columbia.edu/tag/05WS}{Tag 05WS}
(Lemma~10.156.4).

\bibitem{Tamme}
G.~Tamme, \emph{Karoubi's relative Chern character, the rigid syntomic regulator, and the Bloch--Kato exponential map},
Forum Math. Sigma \textbf{2} (2014), e20.
Corollary~5.19 and its proof give the syntomic-to-\'etale compatibility used for $\Spec\OO_E$.
\href{https://arxiv.org/abs/1111.4109v4}{arXiv:1111.4109v4}
(21 July 2014).

\bibitem{TotaroMilnor}
B.~Totaro, \emph{Milnor $K$-theory is the simplest part of algebraic
$K$-theory}, $K$-Theory \textbf{6} (1992), 177--189.
Theorem~1 and Sections~1--2.
\href{https://www.math.ucla.edu/~totaro/papers/public_html/milnor.pdf}
{Author's copy}.

\bibitem{WeibelChern}
C.~Weibel, \emph{\'Etale Chern classes at the prime 2},
in P.~G.~Goerss and J.~F.~Jardine (eds.),
\emph{Algebraic $K$-Theory and Algebraic Topology},
NATO ASI Series C \textbf{407}, Kluwer, 1993, 249--286.
Theorem~3.11(i) gives the line-bundle product formula in $K$-degree three;
Section~5 recalls the indecomposable $K_3$ comparison.
Page locators 10 and 12 refer to the author's version.
\href{https://sites.math.rutgers.edu/~weibel/archive/papers-dir/chernclass.pdf}{Author's paper};
\href{https://doi.org/10.1007/978-94-017-0695-7_14}{doi:10.1007/978-94-017-0695-7\_14}.

\bibitem{Weibel}
C.~A.~Weibel, \emph{The $K$-book: An Introduction to Algebraic $K$-theory}, Graduate Studies in Mathematics \textbf{145}, American Mathematical Society, 2013.
Chapter~III, Definition~5.11 and relations (D1)--(D3), pp.~43--44,
fix the Dennis--Stein convention used here; Chapter~III,
Theorem~6.2.4, p.~50, gives the local-field $K_2$ decomposition.
See also Chapter~IV, Theorem~1.12 and Corollary~1.13, p.~10,
for finite fields, Section~2 for coefficients, and Section~6.4
for filtered colimits; Chapter~V, 3.3, 3.3.2, 3.7.2, 3.12,
4.1, and 6.1 for the cited transfers, base change, projection,
d\'evissage, and localization results; and Chapter~VI,
Theorems~4.1 and~5.2 for norm-residue and the rational
Bloch-group identification. Page locators refer to the author PDFs.
\href{https://sites.math.rutgers.edu/~weibel/Kbook/Kbook.III.pdf}{Author's Chapter~III};
\href{https://sites.math.rutgers.edu/~weibel/Kbook/Kbook.IV.pdf}{Author's Chapter~IV};
\href{https://sites.math.rutgers.edu/~weibel/Kbook/Kbook.V.pdf}{Author's Chapter~V};
\href{https://sites.math.rutgers.edu/~weibel/Kbook/Kbook.VI.pdf}{Author's Chapter~VI}.
\end{thebibliography}
